\section{Introduction}
%\vspace{-1ex}

% \ak{High-level point: only one main point per paragraph! Especially the later paragraphs are covering way too many topics in a single paragraph. I would come up with the first sentence for each paragraph that will describe the full paragraph, and only stick to things that support the first sentence}

% \ak{High-level point: 1) A lot of sentences would be cleaner if split into two sentences, instead of a long run-on sentence 2) Use commas!}

Molecular conformer generation (MCG) is a fundamental task in computational chemistry. The objective is to predict stable low-energy 3D molecular structures, known as conformers. 
Accurate molecular conformations are crucial for various applications that depend on precise spatial and geometric qualities, including drug discovery and protein docking.
For MCG, traditional physics-based methods present a trade-off between speed and accuracy. Quantum mechanical methods are more accurate but computationally slow, while stochastic cheminformatics-based methods like RDKit EKTDG~\citep{rdkit} provide more efficient but less accurate results.
As the difficulty of computing low-energy structures increases with the number of atoms and rotatable bonds in a molecule, there has been interest in developing machine learning (ML) methods to generate accurate conformer predictions efficiently. 

% \begin{figure}[t]
% \centering
%     \includegraphics[width=\textwidth]{figures/figure2.png}
%     %\vspace{-10pt}
%     \caption{\danny{move to supplementary}\textbf{Coarse-grained prediction task.} \textbf{(I)} Fine-grain molecules are split along rotatable bonds that define torsion angles.They are then coarse-grained to reduce the dimensionality and learn a subgraph-level latent distribution. We stress N can differ for each molecule depending on the number of torsion angles. \textbf{(II)} Illustration of the decoder message passing parameterization and learning objective. We learn the optimal distortion, \ie\ distance from the reference molecule. 
%      }
%     \label{fig:novelty}
%     %\vspace{-5pt}
% \end{figure}

% \ak{In this section, describe challenges of prior generative models. To do better at this task, we need more flexibility + a better representation (go into detail, and/or examples). Our CG + attention mechanism gives both (list reasons why). Then we go into the full architecture/method. }

% Existing generative MCG models can be distinguished by two main criteria: their (1) probabilistic modeling architecture and the (2) geometric information to model 3D conformers.
% Most methods either use distance matrices, atom coordinates, or torsional angles, all of which have been applied to various probabilistic modeling techniques, from variational autoencoders to diffusion models.
% While successful, a common limitation of these prior works is that they limit themselves to modeling one source of geometric information. By only working in pure coordinate or angle space, the inductive bias of the equivariance of the problem is not taken full advantage of.
% As a result, we introduce a more flexible molecule latent representation that better integrates various desired geometric modalities.
Existing generative MCG ML models can be broadly classified based on two primary criteria: (1) the choice of the architecture to model the distribution of low-energy conformers and (2) the manner in which they incorporate geometric information to model 3D conformers. A majority of these methods utilize a single geometric representation, such as distance matrices, atom coordinates, or torsional angles, and pair them with various probabilistic modeling techniques, including variational autoencoders (VAEs) and diffusion models. 
A recurring limitation of these prior approaches is that they are restricted to modeling only one specific aspect of geometric information \ie, coordinates or angles. 
By restricting themselves to either a pure coordinate or angle space, they often fail to fully leverage the full geometric information inherent to the problem. Consequently, in this work, we introduce a more flexible molecular latent representation that seamlessly integrates multiple desired geometric modalities, thereby enhancing the overall accuracy and versatility of conformer generation.

% The basis for our work stems from prior 2D fragment-based generative techniques in which breaking the molecule into prominent motifs or substructures resulted in a significant performance boost in both learning optimal representations and generative tasks. 
% This same idea can be applied to 3D space via the coarse-graining of input molecules to condense information into a coarse-grained (CG) representation based on desired coarsening criteria.
% Coarse-graining has shown great success in related tasks such as molecular dynamics and protein modeling, but these related methods have only explored rigid coarsening criteria, such as a distance-based nearest neighbor approach, to ensure all inputs can be represented by a fixed N CG beads. 
% In many applications, not every input can be meaningfully forced into the same granularity, especially small molecules. 
% As a result, we propose an aggregated attention mechanism that enables molecules of any size and any number of CG beads to be represented in a single latent representation. The flexibility the attention provides allows input molecules to be fragmented along torsion angles to model interatomic distances, 3D atom coordinates, and torsion angles directly or implicitly together in an equivariant fashion.

% \ak{Start by introducing coarse-graining in 3D (the main idea), and then talk about overlap with previous 2D tasks. First sentence should always be the motivator for the rest of the paragraph} \danny{its hard to define what CG is without equations or in a small amount so I am trying to use the connection from 2D hence the order}
% %\vspace{-3ex}
% \ak{Can we talk about first sentence which sets up the whole paragraph. Maybe natural dimensional reduction to learn better, etc. Then 2D does it this way, we choose CG to do this in 3D}
% While MCG is a 3D task, we draw inspiration from previous 2D fragment-based generative techniques~\citep{FaST}. In these scenarios, the fragmentation of molecules into prominent substructures led to significant enhancements in representation learning and generative tasks. Extending this concept into 3D space involves a process known as coarse-graining, which condenses molecular information into a coarse-grained (CG) representation based on specific coarsening criteria. \danny{work this last sentence as the last sentence}

To create this flexible representation, we use a process known as coarse-graining to distill molecular information into a simplified, coarse-grained (CG) representation based on specific coarsening criteria. This is analogous to previous 2D fragment-based generative techniques~\citep{FaST}. In these scenarios, the fragmentation of molecules into prominent substructures led to significant enhancements in representation learning and generative tasks.
Coarse-graining has seen success in related ML applications like molecular dynamics and protein modeling~\citep{CG-MD, CGProt-background}. However, these ML approaches have primarily explored rigid coarsening criteria, such as distance-based nearest neighbor methods, which represent all inputs with a fixed granularity or number of beads (\ie\ CG nodes, $N$). 
ML models that use fixed-length coarse-graining often necessitate $N$ be fixed for all input molecules~\citep{CGProtein}. 
This approach may not be suitable for all scenarios, especially when navigating multi-modal datasets of drug-like molecules of varying sizes (multiple $N$), which better reflects real-world conditions.

To address the limitations of fixed-length CG representations, we introduce \textit{Aggregated Attention} for variable-length coarse-graining (see~\fref{fig:variable_legnth}). 
This methodology allows a single latent representation to accommodate molecules with different numbers of fine-grained (FG) atoms and CG beads.
The inherent flexibility of the attention mechanism allows input molecules to be fragmented along torsion angles, enabling the modeling of interatomic distances, 3D atom coordinates, and torsion angles in an equivariant manner, regardless of the molecule's shape or size.
% We provide an example of variable-length coarse-graining in~\fref{fig:variable_legnth}.
% Unlike prior methods that are limited to fixed-length coarse-graining, our approach learns variable-length CG representations. 
% This adaptability improves our capacity to generate more accurate structures, particularly across datasets featuring molecules of varying sizes, which better reflects real-world conditions. 
Through Aggregated Attention, we also harness information from the entire learned representation, unlike preceding approaches that restrict the backmapping (from CG back to FG) to a subset~\citep{CGVAE}. % of channels (top-k channels)~\citep{CGVAE}.
% In contrast to previous methods, which are confined to fixed-length coarse-graining, our approach learns variable-length CG representations. This adaptability enhances our ability to generate more accurate structures, particularly across datasets featuring molecules of widely varying sizes—--a scenario that is much more reflective of real-world conditions. Aggregated Attention makes use of information from the entire learned representation, in contrast to previous methods that limit the backmapping (from CG back to FG) to a subset of channels (top-k channels)~\citep{CGVAE}.
% The adaptability of our learned representations improves the capacity to generate more accurate structures. 
The adaptability of our learned variable-length representations enables more accurate generation.

\begin{figure}[t]
\centering
    \includegraphics[width=0.8\textwidth]{figures/Variable_Length3.png}
    % %\vspace{-10pt}
    % %\vspace{-4ex}
    \caption{\textbf{Learning flexible coarse-grained representations.}  \model is the first model to employ variable-length (on left) coarse-graining. Each input molecule ($n$ fine-grained (FG) atoms) can be represented by a different number of coarse-grained (CG) nodes $N$, thus accommodating diverse molecular sizes. In contrast, prior approaches rely on fixed-length (on right)  coarse-graining, thereby forcing all molecules to possess the same number of CG nodes. Variable-length coarse-graining enhances the model's ability to create better learned representations across molecules of different sizes and geometries. The molecules on the left are coarsened along torsional angles.}
    \label{fig:variable_legnth}
    % %\vspace{-5pt}
    % %\vspace{-3ex}
\end{figure}

Our full innovations are encapsulated in \modelns, an SE(3)-equivariant hierarchical VAE. 
\model aggregates information from fine-grained atomic coordinates to create a flexible subgraph-level representation, improving the accuracy of conformer generation.
% In contrast to previous more general backmapping techniques~\citep{CGVAE} that only utilize a subset of the learned representation, \modelns's novel Aggregated Attention mechanism, while enabling variable-length coarse-graining, leverages information from the entire learned representation. This prevents wasted computation and allows the information from the hierarchical encoder to better inform the backmapping (CG back to FG) process.
Unlike prior MCG methods, \model generates low-energy conformers with the ability to model atomic coordinates (FG and CG), distances, and torsion angles directly via variable-length coarse-graining. 

Our main contributions are as follows:

\begin{itemize}
% %\vspace{-3ex}
\item We present \modelns, a novel conditional hierarchical VAE. \model learns a coarse-grained or subgraph-level latent distribution for SE(3)-equivariant conformer generation.
% Our coarse-graining strategy allows \model to learn optimal torsional angles implicitly via autoregressive generation or explicitly by adding angle-based loss functions or angle prediction.
To our knowledge, this is the first method to use coarse-graining in the context of MCG.% and the first to explore 3D autoregressive techniques for MCG.

\item \model is the first model capable of handling variable-length coarse-to-fine generation using an \textit{Aggregated Attention} strategy. 
\model employs a single flexible variable-length node-level latent representation that can uniquely represent molecules of any size with any number of coarse-grained nodes. 
% Our variable-length flexibility eliminates the need to train separate models for each distinct number of CG nodes (necessary for prior fixed-length CG methods), which would involve hundreds of models for MCG.
% Utilizing variable-length coarse-graining enables us to circumvent the limitation that fixed-length coarse-graining has of having to train separate generative models for each $N$ (100+ for MCG) to represent the same molecular dataset accurately.
Furthermore, variable-length coarse-graining circumvents having to train separate generative models for each number of CG beads to represent the same molecular dataset accurately (100+ models for MCG), which is a limitation of fixed-length methods~\citep{CGProtein}.

% \item We predominantly outperform prior methods on GEOM-QM9 and GEOM-DRUGS evaluation benchmarks \citep{GEOM}. We do so while requiring less training time, 2\% the total training iterations, and 59\% and 92.7\% less data compared to prior state-of-the-art for QM9 and DRUGS, respectively.
\item We predominantly outperform prior methods on GEOM-QM9 and GEOM-DRUGS for RMSD precision and property prediction benchmarks \citep{GEOM}. We also produce a lower overall RMSD distribution across all conformers, achieving this with an order of magnitude less training time compared to prior methods. 
% We note that by using an autoregressive decoder, we require overall less training time and can also leverage 59\% and 92.7\% less data compared to prior state-of-the-art for QM9 and DRUGS, respectively.

% \item We enhance the existing benchmarks for conformer generation by incorporating direct measurements of average and maximum root mean square deviation (RMSD) across the set of generated conformers for each molecule. These new metrics provide a more comprehensive assessment of a model's ability to generate robust conformations, as they evaluate the quality of all generated samples, not just the best ones.
% Here, robust refers to the consistent generation of high-quality conformers.
\item We evaluate \model on multiple downstream applications to assess the chemical and biochemical quality of our generated conformers, including oracle-based protein docking~\citep{tdc} (the affinity of generated conformers to bind to specific protein pockets) under both flexible and rigid conformational energy minimizations. Despite lacking prior knowledge about the protein or the downstream task, \model generates significantly better binding ligand conformers for known protein binding sites when compared to prior MCG methods for both oracle scenarios.

% incorporate a practical application of conformers, oracle-based protein docking~\citep{tdc}, to assess the structural precision and biochemical quality of MCG models.% through biologically relevant metrics. 
% We focus on two distinct oracle function scenarios— flexible and rigid conformational energy minimizations. 
% Leveraging standardized prediction software and data from structure-based drug design (SBDD)~\citep{targetdiff}, we assess the quality of generated conformers based on their affinity to bind to specific protein pockets. 
% We highlight how \modelns, despite lacking prior knowledge about the protein or the downstream task, excels in generating significantly better binding ligand conformers for known protein binding sites when compared to prior MCG methods for both oracle scenarios.

% We integrate a common practical use of conformers, oracle-based protein docking~\citep{tdc}, to compare the structural accuracy of MCG models via biologically meaningful measurements. We isolate two protein docking tasks (with and without explicit energy minimization) and use standardized binding-affinity prediction software and large-scale structure-based drug design datasets to evaluate generated conformers by their ability to bind to target protein locations. We demonstrate how \modelns, without knowing anything about the protein or the task itself, can generate significantly more accurate ligand conformers for known protein pockets with and without protein-dependent conformational energy minimizations.

% \item We use standardized binding-affinity prediction software and a large-scale structure-based drug design datasets for ligand-protein binding affinity-based conformer evaluations. For known proteins with known binding pockets, the more accurate the generated ligand structures are, \ie its conformer, the better the calculated binding affinities will be. As such, we isolate two (with and without explicit energy minimization) ligand-centric tasks by using protein docking as the tool to measure the structural accuracy of drug-like molecules via biologically meaningful measurements (free energy change of binding).
\end{itemize}
%\vspace{-2ex}